\chapter{Συμπεράσματα}
\label{ch:conclusions}

Η παρούσα διπλωματική εργασία ανέπτυξε ένα πλαίσιο προσομοίωσης
link-budget για ένα σενάριο STIN με δύο επίγειους σταθμούς βάσης 5G NR
και έναν δορυφόρο LEO, βασισμένο σε πρότυπα μοντέλα διάδοσης 3GPP
TR~38.901 για το επίγειο σκέλος και σε απλοποιημένο μοντέλο ελεύθερου
χώρου για το δορυφορικό, με single-connectivity επιλογή κόμβου βάσει
μέγιστου SNR.

Το βασικό βήμα πέρα από το αρχικό στατικό σενάριο ήταν η προσομοίωση
πολλαπλών διαδοχικών μεταδόσεων στη διάρκεια μίας διέλευσης δορυφόρου
LEO, με κίνηση υπολογισμένη από την τροχιακή μηχανική αντί για αυθαίρετη
παραδοχή. Η επέκταση αυτή, απούσα από την αρχική έκδοση της
προσομοίωσης, ανέδειξε τα πιο χρήσιμα ευρήματα της εργασίας: μια
ξεκάθαρη ισορροπία μεταξύ ρυθμού και ενεργειακής απόδοσης ανάμεσα στις
δύο ζεύξεις, και δύο συγκεκριμένες αδυναμίες του κανόνα επιλογής κόμβου
που κανένα μεμονωμένο στατικό σενάριο δεν θα μπορούσε να φανερώσει. Η
πρώτη -- ασυσχέτιστο shadow fading που παρήγε τεχνητά συχνά handovers --
αντιμετωπίστηκε ήδη μέσα στην ίδια εργασία, με την υιοθέτηση ενός χωρικά
συσχετισμένου μοντέλου σκίασης· η δεύτερη, η απουσία ρητής διαχείρισης
outage, παραμένει ανοιχτή και ορίζει την επόμενη προτεραιότητα.

Ένα πρώτο μοντέλο μηχανικής μάθησης επιβεβαίωσε ότι η ροή δεδομένων
MATLAB $\rightarrow$ CSV $\rightarrow$ Python λειτουργεί σωστά από άκρη
σε άκρη, δεν αποτελεί όμως ακόμα μια πραγματικά δύσκολη πρόβλεψη -- κάτι
που παραμένει στόχος του επόμενου σταδίου.

Το γενικότερο συμπέρασμα είναι ότι η μετάβαση από μία στατική
προσομοίωση σε χρονική διάσταση δεν προσθέτει απλώς όγκο δεδομένων, αλλά
φέρνει στην επιφάνεια πραγματικά ζητήματα σχεδίασης, όπως η διαχείριση
outage και η υστέρηση (hysteresis) στην απόφαση handover, τα οποία
παραμένουν αόρατα σε μια στατική ανάλυση. Τα ζητήματα αυτά, μαζί με τη
μετάβαση σε πραγματική πολυσυνδεσιμότητα και σε πιο απαιτητικούς στόχους
μηχανικής μάθησης, καθορίζουν το επόμενο στάδιο της εργασίας, όπως
περιγράφηκε στο Κεφάλαιο~\ref{ch:discussion}.
